% !TEX root = ../main.tex
% ---------------------------------------------------------------------
%  B5 -- Endliche Automaten (Skript Kap. 5, Anhang C.12; AS 4; MEP HS24)
% ---------------------------------------------------------------------
\vlsec{cB5}{B5}{Endliche Automaten}
\label{sec:automaten}

\SW{Zustand} Situation, in der sich das System befindet (Init,
Automatik, Hand, Störung). Immer \emph{genau ein} aktiver Zustand.
\SW{Transition} Bedingung für den Wechsel: Sensor, Tastendruck (Flanke!),
abgelaufener Timer, Kombination davon.
\SW{Diagramm} Zustände als Kreis, Transitionen als Pfeil mit Bedingung,
definierter Startzustand (Init). Steuerung, keine Regelung.

\subsection{Zwei Implementierungen in TwinCAT}
\begin{ztab}{@{}lX@{}}
UML-Statechart & grafisch, pro Zustand \T{ENTRY} (1x beim Eintreten),
  \T{DO} (zyklisch, solange aktiv), \T{EXIT} (1x beim Verlassen);
  Transitionen mit Bedingung. Gut für Kommunikation. \\
ST: ENUM + CASE & Zustandsvariable vom ENUM-Typ, \R{CASE} springt in den
  Codeblock, Wechsel durch Zuweisung. Prüfungsrelevant. \\
\end{ztab}

\begin{rezept}[Automat in ST -- Grundschema]
\label{rez:automat}
\begin{enumerate}
\item ENUM mit allen Zuständen unter DUTs (Startzustand als Initialwert).
\item Eingänge einlesen und vorverarbeiten (Flanken, Timer \emph{vor} dem
  \R{CASE} aufrufen).
\item \R{CASE eState OF}: pro Zustand Aktionen, am Ende Transition prüfen
  und \R{eState := NEU;} setzen.
\item Ausgänge \emph{nach} dem \R{CASE} aus dem Zustand ableiten
  (\R{bMotor := eState = E\_RUN;}) \ra kein Ausgang bleibt hängen.
\item Logik in einen \R{FUNCTION\_BLOCK}, in \R{MAIN} instanziieren
  \ra mehrfach verwendbar.
\end{enumerate}
\end{rezept}
\begin{falle}[Typische Fehler]
Taster ohne Flanke abfragen \ra Zustand rast in einem Zyklus durch
mehrere Zustände. Timer nur im \R{CASE}-Zweig aufrufen \ra läuft nicht
weiter. Störung nur in einem Zustand prüfen \ra wird übersehen (globale
Transition \emph{vor} dem \R{CASE}).
\end{falle}

\subsection{MEP HS24: Förderband}
\label{bsp:foerderband}
\begin{center}
\begin{tikzpicture}[>=Stealth, node distance=8mm and 13mm,
  every state/.style={minimum size=9mm, inner sep=1pt, font=\scriptsize\bfseries,
    draw=cB5, fill=cB5!8}, every edge/.style={draw, ->}, font=\tiny]
\node[state, initial, initial text=] (I) {INIT};
\node[state, right=of I] (R) {READY};
\node[state, right=of R] (U) {RUNNING};
\node[state, below=of U] (S) {STOPPED};
\node[state, below=of R] (F) {FAULT};
\path (I) edge node[above] {sofort} (R)
      (R) edge node[above] {Start} (U)
      (U) edge node[right] {Stop} (S)
      (S) edge node[above, sloped] {Stop los} (R)
      (F) edge node[left] {Reset} (R);
\node[font=\tiny, below=1mm of F] {von jedem Zustand: FaultSignal \ra FAULT};
\end{tikzpicture}
\end{center}
\begin{lstlisting}
TYPE EN_FOERDERBAND_ZUSTAND :
  (INIT, READY, RUNNING, STOPPED, FAULT);
END_TYPE

FUNCTION_BLOCK FB_FOERDERBAND_LOGIK
VAR_INPUT
  StartButton, StopButton, FaultSignal, ResetButton : BOOL;
END_VAR
VAR_OUTPUT ConveyorOn, FaultIndicator : BOOL; END_VAR
VAR
  eZustand : EN_FOERDERBAND_ZUSTAND := INIT;
  rtStart, rtReset : R_TRIG;
END_VAR

rtStart(CLK := StartButton);
rtReset(CLK := ResetButton);
IF FaultSignal THEN eZustand := FAULT; END_IF  // gilt überall

CASE eZustand OF
  INIT:    eZustand := READY;
  READY:   IF rtStart.Q THEN eZustand := RUNNING; END_IF
  RUNNING: IF StopButton THEN eZustand := STOPPED; END_IF
  STOPPED: IF NOT StopButton THEN eZustand := READY; END_IF
  FAULT:   IF rtReset.Q AND NOT FaultSignal THEN
             eZustand := READY;
           END_IF
END_CASE

ConveyorOn     := (eZustand = RUNNING);
FaultIndicator := (eZustand = FAULT);
\end{lstlisting}
\begin{lstlisting}
PROGRAM MAIN
VAR
  bStart  AT %I* : BOOL;   bStop  AT %I* : BOOL;
  bFault  AT %I* : BOOL;   bReset AT %I* : BOOL;
  bMotor  AT %Q* : BOOL;   bLampe AT %Q* : BOOL;
  fbBand1 : FB_FOERDERBAND_LOGIK;
END_VAR
fbBand1(StartButton := bStart, StopButton := bStop,
        FaultSignal := bFault, ResetButton := bReset,
        ConveyorOn => bMotor, FaultIndicator => bLampe);
\end{lstlisting}
\begin{formel}[(e) Zweites Förderband]
\label{falle:instanz}
Zweite \T{Instanz} \R{fbBand2 : FB\_FOERDERBAND\_LOGIK;} plus eigene
E/A-Variablen (\R{AT \%I*}/\R{AT \%Q*}) deklarieren, verknüpfen und
\R{fbBand2(...)} aufrufen. Der FB-Code bleibt unverändert -- jede Instanz
hat ihren eigenen \R{eZustand}. Genau darum gehört die Logik in einen FB
und nicht ins Programm.
\end{formel}

\subsection{MEP HS24: Ampel für eine Kreuzung}
\label{bsp:ampel}
\begin{lstlisting}
TYPE EN_FARBEN        : (ROT, GELB, GRUEN);      END_TYPE
TYPE EN_BETRIEBSMODUS : (AUTOMATIK, MANUELL);    END_TYPE
TYPE EN_PHASE : (A_GRUEN, A_GELB, B_GRUEN, B_GELB); END_TYPE

VAR_GLOBAL   // GVL_IO
  bA_Rot AT %Q* : BOOL; bA_Gelb AT %Q* : BOOL; bA_Gruen AT %Q* : BOOL;
  bB_Rot AT %Q* : BOOL; bB_Gelb AT %Q* : BOOL; bB_Gruen AT %Q* : BOOL;
  bTA_Rot AT %I* : BOOL; bTA_Gelb AT %I* : BOOL; bTA_Gruen AT %I* : BOOL;
  bTB_Rot AT %I* : BOOL; bTB_Gelb AT %I* : BOOL; bTB_Gruen AT %I* : BOOL;
  bHand   AT %I* : BOOL;              // Wahlschalter Handbetrieb
END_VAR

PROGRAM MAIN
VAR
  eModus  : EN_BETRIEBSMODUS := AUTOMATIK;
  arAmpel : ARRAY[1..2] OF EN_FARBEN := [ROT, ROT];
  ePhase  : EN_PHASE := A_GRUEN;
  tonGruen, tonGelb : TON;
END_VAR

IF GVL_IO.bHand THEN eModus := MANUELL; ELSE eModus := AUTOMATIK; END_IF

CASE eModus OF
AUTOMATIK:
  tonGruen(IN := (ePhase = A_GRUEN) OR (ePhase = B_GRUEN), PT := T#10S);
  tonGelb (IN := (ePhase = A_GELB)  OR (ePhase = B_GELB),  PT := T#3S);
  CASE ePhase OF
    A_GRUEN: arAmpel[1] := GRUEN; arAmpel[2] := ROT;
             IF tonGruen.Q THEN ePhase := A_GELB; END_IF
    A_GELB:  arAmpel[1] := GELB;  arAmpel[2] := ROT;
             IF tonGelb.Q THEN ePhase := B_GRUEN; END_IF
    B_GRUEN: arAmpel[1] := ROT;   arAmpel[2] := GRUEN;
             IF tonGruen.Q THEN ePhase := B_GELB; END_IF
    B_GELB:  arAmpel[1] := ROT;   arAmpel[2] := GELB;
             IF tonGelb.Q THEN ePhase := A_GRUEN; END_IF
  END_CASE
MANUELL:   // Grün nur, wenn die andere Ampel nicht grün ist
  IF GVL_IO.bTA_Rot THEN arAmpel[1] := ROT;
  ELSIF GVL_IO.bTA_Gelb THEN arAmpel[1] := GELB;
  ELSIF GVL_IO.bTA_Gruen AND arAmpel[2] <> GRUEN THEN arAmpel[1] := GRUEN;
  END_IF
  IF GVL_IO.bTB_Rot THEN arAmpel[2] := ROT;
  ELSIF GVL_IO.bTB_Gelb THEN arAmpel[2] := GELB;
  ELSIF GVL_IO.bTB_Gruen AND arAmpel[1] <> GRUEN THEN arAmpel[2] := GRUEN;
  END_IF
END_CASE

GVL_IO.bA_Rot := arAmpel[1] = ROT;  GVL_IO.bA_Gelb := arAmpel[1] = GELB;
GVL_IO.bA_Gruen := arAmpel[1] = GRUEN;
GVL_IO.bB_Rot := arAmpel[2] = ROT;  GVL_IO.bB_Gelb := arAmpel[2] = GELB;
GVL_IO.bB_Gruen := arAmpel[2] = GRUEN;
\end{lstlisting}
\begin{falle}[Ampel: was man begründen sollte]
Zwei Timer reichen, weil zwischen zwei Grünphasen immer eine Gelbphase
liegt (\R{IN} fällt dazwischen ab \ra Reset). Rot einer Ampel dauert hier
Grün $+$ Gelb der anderen $=13$\,s; exakt 10\,s Rot ginge nur mit
kürzerer Grünphase der Gegenrichtung. „Nie beide grün“ ist in Automatik
durch die Phasen garantiert, im Handbetrieb durch die Verriegelung.
\end{falle}

\subsection{Weitere Automaten aus AS 4}
\SW{Safe} Init (Code-Array füllen) \ra Close (Zähler/Lampen reset) \ra
CheckDigit (nur bei \R{risingEdge}, 4 Tasten, \R{correctCounter},
\R{currentIndex}) \ra Open (grün, \R{lightDuration} 5\,s) bzw.\ WrongCode
(rot, 5\,s) \ra Close. Zusatz: Doppelklick blau bricht ab; Code im Zustand
Open mit beiden Tasten neu programmieren.
\SW{Ampel-FB} Eingänge \R{clear\_road}, \R{block\_road},
\R{disable\_system}; Init \ra Red \ra(clear) Red\&Yellow \ra(2\,s) Green
\ra(block) Yellow \ra(5\,s) Red; Disabled: gelb blinkt 1\,s, zurück zu Red,
wenn \R{disable} weg. Als FB \ra Kreuzung $=$ 4 Instanzen, Phase A
(Ampel 1+3, 5\,s), Phase B (2+4, 8\,s), dazwischen sicherer Wechsel.
